\begin{table}[!ht]\centering
\caption{Descriptive statistics}\label{tab:desc}
\begin{threeparttable}\footnotesize\setlength{\tabcolsep}{3pt}
\begin{tabular}{lccccc}\toprule
Variable & N & Mean & S.D. & Min & Max \\ \midrule
\multicolumn{6}{l}{\emph{Panel A. Occupation-year cells, 2014--2020}} \\
Log mean gross hourly pay & 2,546 & 2.675 & 0.367 & 1.908 & 4.052 \\
Mean gross hourly pay (\pounds) & 2,546 & 15.61 & -- & 6.74 & 57.51 \\
Employee jobs (thousands) & 2,303 & 78.7 & -- & 4 & 1178 \\
Mean paid hours per week & 2,551 & 35.1 & -- & -- & -- \\
\addlinespace
\multicolumn{6}{l}{\emph{Panel B. Occupation-level attributes, fixed over time}} \\
Automation probability & 366 & 0.446 & 0.145 & 0.181 & 0.728 \\
Share of jobs below \pounds 7.20 in 2015 & 347 & 0.086 & 0.145 & 0.000 & -- \\
\bottomrule\end{tabular}
\begin{tablenotes}\footnotesize
\item \emph{Notes:} Pay is nominal. Jobs and hours are the published occupation totals and means from ASHE Tables 14.5a and 14.9a; the jobs count is published for fewer cells because the ONS suppresses small counts. The automation probability is defined on the SOC 2010 four-digit classification. The share of jobs below \pounds 7.20, the National Living Wage rate of April 2016, is interpolated from the published percentiles of the 2015 release by equation (B1); its tenth and ninetieth percentiles across occupations are 0.000 and 0.286, and it correlates 0.60 with the automation probability.
\item \emph{Source:} ASHE Table 14, 2014--2020 releases, revised editions; ONS (2019) probability of automation.
\end{tablenotes}\end{threeparttable}\end{table}